\icode{E-MI-SI-L3-261109-001}
如图所示，轻弹簧一端连接小球 \(A\)，另一端悬挂在天花板上。轻绳 1 连接小球 \(A\)、\(B\)，轻绳 2 一端连接小球 \(B\)，另一端固定在竖直墙壁上，两小球均处于静止状态。\(A\)、\(B\) 的质量分别为 \(M\)、\(m\)，弹簧弹力为 \(F\)，弹簧与竖直方向、轻绳 1 与水平方向的夹角均为 \(\theta =\) \,\qty{30}{\degree}，轻绳 2 沿水平方向，重力加速度为 \(g\)。下列说法正确的是\paren[]
	
\begin{tikzpicture}[font=\scriptsize]
	\draw[magenta] (0,0)--(-60:.5)coordinate(e) (-60:1.5)--(-60:2)coordinate(A)--++(-30:1.2)coordinate(B)--++(1,0);
	\draw[dashed,cyan] (0,0)--++(-90:1.2) (B)--++(-1.2,0);
	\draw[gray] (0,-.3) arc(-90:-60:.3) ($(B)+(-.4,0)$) arc(180:150:.4);
	\draw[thick,decoration={aspect=.6,segment length=3,amplitude=.8mm,coil},decorate] (e)--(-60:1.5);
	\foreach \x in {A,B}
		\draw[fill=white] (\x)node{ \(\x\)} circle(.2);
	\foreach \x in {-.5,-.3,...,.5}
		\draw[gray] (\x,0)--++(.2,.2);
	\foreach \y in {-2.8,-2.6,...,-1.8}
		\draw[gray] (3.06,\y)--++(.2,.2);
	\draw[thick] (-.5,0)--(.5,0) (3.06,-2.8)--(3.06,-1.8);
	\node at(1.55,-1.8){ \(1\)};
	\node at(2.55,-2.1){ \(2\)};
	\node at(.12,-.6){ \(\theta\)};
	\node at(1.4,-2.15){ \(\theta\)};
\end{tikzpicture}
\begin{choices}
	\item \(M = \sqrt{3}m\)
	\item \(M = m\)
	\item \(F = 2\sqrt{3}mg\)
	\item \(F = \sqrt{3}mg\)
\end{choices}
\begin{answer}
    C\qquad 提示：
\end{answer}